% Section: Budget-Paced LLM Routing
% Include from paper/sections/results.tex

\subsection{Budget-Paced Routing (Figure~\ref{fig:budget_pacing})}
\label{sec:budget_pacing}

\paragraph{Motivation.}
Production LLM deployments operate under cost constraints:
an engineering team may route millions of requests per day across
models spanning a $\bpPriceRangeX\times$ price range
(e.g., \bpFixedLlamaCost/req for Llama-3.1-8B vs.\
\bpFixedGeminiCost/req for Gemini-2.5-Pro)
and must keep average spend within an API budget.
A fixed model gives the operator exactly one operating point
in quality--cost space: choose Llama for low cost and
lower quality ($\bpFixedLlamaReward$), Gemini for high quality
at high cost ($\bpFixedGeminiReward$), or Mistral as a
compromise ($\bpFixedMistralReward$ at \bpFixedMistralCost/req).
None of these can be adjusted without switching models entirely.
We ask: \emph{can a budget-paced router provide a continuous
``dial'' between these fixed points, letting the operator
specify a dollar budget and automatically achieve it?}

\paragraph{Algorithm.}
We adapt the Primal-Dual decomposition of the \emph{Contextual Bandit with Knapsacks}
(CBwK) framework~\cite{agrawal2014bandits} to a per-request rate constraint suited for open-ended deployment (see Section~\ref{sec:problem} for the theoretical status of this adaptation).
At each step $t$ the router selects model $a_t$ by maximising a
budget-augmented utility:
\begin{equation}
    a_t = \arg\max_{a \in \mathcal{A}_t} \Bigl[
        \underbrace{\hat\theta_a^\top x_t + \alpha \sqrt{x_t^\top A_a^{-1} x_t}}_{\text{LinUCB score~\cite{li2010contextual}}}
        \;-\; \underbrace{\lambda_s \,\tilde{c}_a}_{\text{static penalty}}
        \;-\; \underbrace{\lambda_t \,\tilde{c}_a}_{\text{dual penalty}}
    \Bigr]
    \label{eq:budget_ucb}
\end{equation}
where $\mathcal{A}_t \subseteq \mathcal{A}$ is the set of models
surviving a hard cost ceiling
$p_a \leq p_{\max}/(1 + \kappa\lambda_t)$,
$\tilde{c}_a \in [0,1]$ is the log-normalized unit cost,
and $\lambda_t$ is a Lagrange multiplier updated after observing the
realised cost $c_t$.

The dual update uses an \emph{EMA-smoothed} cost signal rather than
the raw per-request cost:
\begin{align}
    \bar{c}_t &= (1 - \alpha_{\mathrm{ema}})\,\bar{c}_{t-1}
                 + \alpha_{\mathrm{ema}}\,c_t
    \label{eq:cost_ema} \\
    \lambda_{t+1} &= \min\!\Bigl(\bar\lambda,\;
    \max\bigl(0,\;\lambda_t + \eta\,(\bar{c}_t / B - 1)\bigr)\Bigr)
    \label{eq:dual_update}
\end{align}
with budget target $B$ (the user's desired average cost per request),
EMA smoothing factor $\alpha_{\mathrm{ema}}{=}0.05$,
step size $\eta{=}0.05$, and cap $\bar\lambda{=}5$.
Using $\bar{c}_t$ instead of $c_t$ is critical when model costs span
orders of magnitude: a single Gemini request at \bpFixedGeminiCost\
would otherwise spike $c_t/B$ by up to $\bpPriceRangeX\times$ the target,
producing sawtooth oscillations in $\lambda_t$ that never converge.
The EMA smooths this variance, producing a stable gradient signal
whose magnitude reflects the \emph{average} cost regime.
The normalisation $\bar{c}_t / B$ makes $\eta$ portfolio-independent:
the same step size produces equivalent $\lambda_t$ trajectories regardless
of whether the cheapest model costs \$0.00003 or \$0.03.

The \textsc{BudgetPacer} operates in \emph{adaptive mode}, combining
two enforcement mechanisms:
\begin{itemize}[nosep,leftmargin=*]
  \item \textbf{Hard ceiling} (safety net):
    When $\lambda_t > 0$, the candidate set $\mathcal{A}_t$ excludes
    models whose blended unit price exceeds the dynamic ceiling,
    preventing catastrophic overspend on any single request.
  \item \textbf{Soft penalty} (optimizer):
    The dual penalty $\lambda_t\tilde{c}_a$ biases the UCB score toward
    cheaper models proportionally to their cost, allowing context-dependent
    routing within the surviving candidate set.
\end{itemize}

\noindent
This two-layer design mirrors production rate-limiting architectures:
the hard ceiling is a circuit breaker that guarantees worst-case
containment; the soft penalty is a fine-grained optimizer that
minimizes quality loss within the budget envelope.

\paragraph{Setup.}
We evaluate on the $K{=}3$ portfolio
(Llama-3.1-8B / Mistral-Large / Gemini-2.5-Pro)
under stationary conditions.
The router learns online on the validation split
($n{=}1{,}785$; 32~NLP benchmarks) and is evaluated on a held-out
test split ($n{=}1{,}824$) across 20~seeds.
Following standard bandit regret evaluation protocol, the router
continues learning during the test phase: each prompt is routed
\emph{before} its reward is observed (no look-ahead), and the
bandit parameters update after every decision, matching
the production deployment where the router never stops learning.
Warmup priors are trained on a disjoint training split
($n_{\mathrm{eff}}{\approx}\bpNeff$, $\alpha{=}\bpAlpha$, PCA-25,
$\gamma{=}\bpGamma$; selected via joint epsilon-constraint tuning,
Section~\ref{sec:hparam_sweep}).

Three reference points anchor the evaluation:
the three fixed single-model baselines, each of which always routes
to one model regardless of prompt or budget.
These represent the operating points available to a team that does
not use a router.
The BudgetPacer is evaluated at seven log-spaced targets, ranging
from the cheapest model's empirical mean cost
(\bpCheapestModelCost/req) to the most expensive
(\bpExpensiveModelCost/req).

\paragraph{The router as a quality--cost dial
(Figure~\ref{fig:budget_pacing}a).}
Each fixed model occupies a single point in quality--cost space:
Llama-3.1-8B at (\bpFixedLlamaCost, $\bpFixedLlamaReward$),
Mistral-Large at (\bpFixedMistralCost, $\bpFixedMistralReward$),
and Gemini-2.5-Pro at (\bpFixedGeminiCost, $\bpFixedGeminiReward$).
An operator choosing among these three models faces a discrete
menu with no intermediate options.

The ParetoBandit router, equipped with the BudgetPacer, traces a
\emph{continuous curve} through this space by dynamically mixing
models based on prompt context and budget pressure.
The curve reveals operating points that no single model can reach.
At a budget of \bpAnnotBudget/req, the router achieves
\bpAnnotQualPct\% of Gemini-2.5-Pro quality at only
\bpAnnotCostPct\% of its cost---a region of the quality--cost
tradeoff space that is structurally inaccessible to any fixed
model choice.
This ``intermediate'' quality comes from the router's ability to
blend: at this operating point, it routes \bpAnnotLlamaFrac\% of
traffic through Llama and \bpAnnotMistralFrac\% through Mistral
(Figure~\ref{fig:budget_pacing}c), exploiting each model's
comparative advantage on different prompt types.

At the unconstrained end (budget $\geq$ \bpExpensiveModelCost/req),
the dual variable decays to $\lambda_t \approx 0$ and the pacer
gracefully becomes a no-op: the router achieves
$\bpPacerLoosestReward$, matching the quality ceiling.
This quality ceiling is itself $\bpCeilingOraclePct\%$ of the per-prompt oracle
($\bpOracleReward$, the reward obtained by always selecting the
highest-quality model for each prompt regardless of cost).
The remaining gap reflects the inherent exploration--exploitation
cost of online learning.

\paragraph{Budget compliance
(Figure~\ref{fig:budget_pacing}b).}
The pacer's primary claim is that it \emph{hits the target}.
Figure~\ref{fig:budget_pacing}b plots realised cost against
budget target; points on the 45-degree diagonal represent
perfect compliance.
For the binding targets, budget utilisation ranges from
\bpUtilBindingLow\ to \bpUtilBindingHigh---effectively
perfect compliance.
The loosest target ($B{=}$\bpExpensiveModelCost/req) shows
under-utilisation (\bpLoosestUtil) because the
router's unconstrained mean cost is already below
this target; the dual variable correctly decays to
$\lambda_t \approx 0$ rather than artificially inflating cost.
This asymmetry is by design: the pacer constrains spend
\emph{from above} but never forces the router to spend more than its
quality-optimal policy would.

\paragraph{Model allocation dynamics
(Figure~\ref{fig:budget_pacing}c).}
The stacked bars reveal the mechanism behind the quality--cost curve.
At the tightest budget (\bpCheapestModelCost/req), the hard ceiling
excludes Gemini entirely and the soft penalty suppresses Mistral,
producing a nearly pure Llama allocation ($>$99.9\%).
As budget increases, Mistral enters the mix---at
\bpAnnotBudget/req the router achieves a \bpAnnotLlamaFrac/\bpAnnotMistralFrac\
Llama/Mistral blend that exploits Mistral's favourable
quality-to-cost ratio on prompts where it matches or exceeds
Llama's quality.
At moderate budgets (\bpPacerFourthCost/req), Mistral dominates
at 73\%, with Llama handling the easy prompts and Gemini reserved
for the hardest ones (2\%).
At the loosest budget, Gemini receives 69\% of traffic and the
router approximates the unconstrained quality-optimal policy.

This smooth transition is driven by the dual variable $\lambda_t$:
as the budget target increases, $\lambda_t$ decreases, relaxing
the cost penalty and allowing progressively more expensive models
into the mix.
The transition is not a simple threshold switch---it is
context-dependent: on any given prompt, the router may still
select a cheap model even at a loose budget if the contextual
LinUCB score indicates the expensive model offers negligible
quality uplift for that particular prompt.

\paragraph{Cost at scale.}
The practical value of the continuous quality--cost frontier
compounds with traffic volume.
Consider a deployment routing 100{,}000 requests/day.
Always calling Gemini-2.5-Pro costs
$\bpFixedGeminiCost \times 100{,}000 \approx \$1{,}500$/day.
The router at a moderate budget
($B{=}\bpPacerFifthCost$/req) achieves $\bpPacerFifthReward$
quality---\bpPacerFifthQualPctGemini\% of Gemini---while spending
only \bpPacerFifthCost/req, a
\bpPacerFifthSavingPctGemini\% cost reduction.
At the annotated operating point
(\bpAnnotQualPct\% of Gemini quality), cost drops to
\bpAnnotCost/req---a
\bpAnnotSavingPctGemini\% reduction---while quality remains
well above Llama's baseline.

\paragraph{Production implications.}
\begin{enumerate}[leftmargin=*,noitemsep]
  \item \textbf{Interpretable interface.}
    The practitioner specifies ``keep average cost below
    \$X/request''---a quantity visible on any billing dashboard.
    No offline tuning, no dimensionless penalty weights,
    no re-calibration when pricing changes.

  \item \textbf{Continuous Pareto coverage.}
    Fixed models offer discrete operating points;
    the router fills the gaps between them.
    This is especially valuable in the low-cost regime where
    the router achieves quality well above the cheapest model
    by selectively routing a small fraction of
    traffic to more capable models.

  \item \textbf{Graceful degradation.}
    When the budget is generous, the pacer does not distort routing:
    $\lambda_t \to 0$ and the quality-optimal policy is recovered.
    A team can set a conservative budget ceiling
    (e.g., $2\times$ expected spend) as a safety net without
    sacrificing quality in normal operation.

  \item \textbf{Budget compliance as an SLA.}
    The \bpUtilBindingLow--\bpUtilBindingHigh\ utilisation range
    at binding targets means the pacer can underpin a cost SLA:
    ``average spend will not exceed $1.10 \times B$ over any
    rolling 1{,}000-request window.''
    A fixed model trivially satisfies a budget SLA, but only at
    one quality level---the router satisfies it at \emph{any}
    quality level the operator chooses.
\end{enumerate}
